\subsubsection{যোগাযোগের অর্থ প্রকাশ: \texttt{<address>}}
নথি বা লেখাটির দায়িত্বশীল ব্যক্তি/প্রতিষ্ঠানের যোগাযোগের তথ্য বোঝাতে \texttt{<address>} ব্যবহার করা যায়। এতে শুধু দেখতে বাঁকা করা নয়, তথ্যের ভূমিকাও প্রকাশ পায়। যেকোনো স্থানের ঠিকানা বা ভ্রমণবিবরণ এই ট্যাগে রাখা যথাযথ নয়।
\par\noindent\begin{minipage}[t]{.50\linewidth}\vspace{0pt}\small
\textbf{\color{chFourTeal}কোড: পেজের যোগাযোগ}\par
\begin{verbatim}
<address>
  Science Club Office<br>
  Room 12, Main Building
</address>
\end{verbatim}
\end{minipage}\hfill
\begin{minipage}[t]{.46\linewidth}\vspace{0pt}\small
\textbf{\color{chFourTeal}ফল ও ব্যাখ্যা}\par
\textit{Science Club Office}\par
\textit{Room 12, Main Building}\par
স্বাভাবিক ব্রাউজার রীতিতে বাঁকা দেখা যেতে পারে। \texttt{br} দ্বিতীয় লাইন শুরু করে; CSS দিয়ে রীতি বদলালেও যোগাযোগের অর্থ থাকে।
\end{minipage}\par
